\section{Discussion}
Data augmentation is a necessary method for achieving state-of-the-art performance \cite{simard2003best,krizhevsky2012imagenet,devries2017dataset,zhang2017mixup,girshick2018detectron,park2019specaugment}. Learned data augmentation strategies have helped automate the design of such strategies and likewise achieved state-of-the-art results \cite{cubuk2018autoaugment,lim2019fast,ho2019population,zoph2019learning}.
In this work, we demonstrated that previous methods of learned augmentation suffers from systematic drawbacks. Namely, not tailoring the number of distortions and the distortion magnitude to the dataset size nor the model size leads to sub-optimal performance.
To remedy this situation, we propose a  simple parameterization for targeting augmentation to particular model and dataset sizes. We demonstrate that RandAugment is competitive with or outperforms previous approaches \cite{cubuk2018autoaugment,lim2019fast,ho2019population,zoph2019learning} on CIFAR-10/100, SVHN, ImageNet and COCO without a separate search for data augmentation policies.

In previous work, scaling learned data augmentation to larger dataset and models have been a notable obstacle. For example, AutoAugment and Fast AutoAugment could only be optimized for small models on reduced subsets of data \cite{cubuk2018autoaugment,lim2019fast}; population based augmentation was not reported for large-scale problems \cite{ho2019population}. The proposed method scales quite well to datasets such as ImageNet and COCO while incurring minimal computational cost (e.g. 2 hyper-parameters), but notable predictive performance gains. An open question remains how this method may improve model robustness \cite{lopes2019improving, yin2019afourier, recht2018cifar} or semi-supervised learning  \cite{xie2019unsupervised}. Future work will study how this method applies to other machine learning domains, where data augmentation is known to improve predictive performance, such as image segmentation \cite{chen2017deeplab}, 3-D perception \cite{ngiam2019starnet}, speech recognition \cite{hinton2012deep} or audio recognition \cite{hershey2017cnn}. In particular, we wish to better understand if or when datasets or tasks may require a separate search phase to achieve optimal performance. Finally, an open question remains how one may tailor the set of transformations to a given tasks in order to further improve the predictive performance of a given model.